% Figure: net lift of a fixed-volume hot-air balloon as a function of the
% internal air temperature, at fixed ambient conditions. Lift rises as the
% inside air thins with heating. No em-dashes.
\begin{figure}[htbp]
\centering
\begin{tikzpicture}
\begin{axis}[
  width=11cm, height=7cm,
  xlabel={internal air temperature $T_i$ (K)},
  ylabel={net lift (N)},
  xmin=290, xmax=420, ymin=-500, ymax=3500,
  axis lines=left,
  grid=major, grid style={enggray!20},
  tick label style={font=\scriptsize},
  label style={font=\small},
  legend style={font=\scriptsize},
  legend pos=north west,
]
% Envelope V = 2200 m^3, ambient rho_a = 1.225 at T_a = 288 K.
% Internal density rho_i = rho_a * T_a / T_i (constant pressure).
% Net lift L = (rho_a - rho_i) g V - W_fixed, W_fixed = 3500 N (envelope+basket+crew).
% Plot the buoyant-minus-fixed lift.
\addplot[engorange, very thick, domain=290:420, samples=60]
  ({x}, {(1.225 - 1.225*288/x)*9.81*2200 - 3500});
\addlegendentry{net lift $L(T_i)$}
% zero-lift line
\draw[engred, dashed, thin] (axis cs:290,0) -- (axis cs:420,0);
\node[engred, font=\scriptsize, anchor=south west] at (axis cs:292,0) {launch threshold};
% ambient temperature marker
\draw[enggray, dashed, thin] (axis cs:288,-500) -- (axis cs:288,3500);
\node[enggray, font=\scriptsize, anchor=south, rotate=90] at (axis cs:295,1600) {ambient $T_a$};
\end{axis}
\end{tikzpicture}
\caption[Hot-air balloon net lift versus internal temperature]{Net lift
of a fixed-volume hot-air balloon against the temperature of the air
inside it, at a fixed ambient state. Heating the trapped air at constant
pressure thins it, $\rho_i=\rho_a\,T_a/T_i$, so the envelope displaces the
same volume of denser outside air and the buoyant force grows. The net
lift is the buoyancy less the fixed weight of envelope, basket, and crew;
the balloon leaves the ground once the curve crosses the launch threshold,
here near $370\,\mathrm{K}$ (about $100\,^{\circ}\mathrm{C}$). Pushing the
temperature higher buys more lift but is bounded by the fabric's rating,
which is why a hot-air balloon is flown near, not far above, that
threshold. Computed for a $2200\,\mathrm{m^3}$ envelope at
$\rho_a=1.225\,\mathrm{kg/m^3}$, $T_a=288\,\mathrm{K}$.}
\label{fig:vol03ch10:hot-air-balloon}
\end{figure}
